\documentclass[10pt,a4paper]{amsart}
\usepackage[margin=19mm]{geometry}
\usepackage{amsmath,amssymb,mathtools}
\usepackage[scr=rsfs,cal=euler]{mathalfa}
\usepackage{xfrac}
\usepackage[unicode,hidelinks,bookmarks=false]{hyperref}
\newcommand{\ie}{i.e.\ }
\newcommand{\cf}{cf.\ }
\newcommand{\resp}{respectively\ }
\newcommand{\Subsection}{\subsection}
\providecommand{\nobreakdash}{\nobreak\hbox{-}}
\setlength{\emergencystretch}{2em}
\multlinegap0pt
\usepackage{amssymb}
\usepackage{xcolor}
\newtheorem{lemma}{Lemma}
\newtheorem{proposition}{Proposition}
\title{Balanced Bismut torsion-parallel fourfolds with constant holomorphic sectional curvature}
\author{Qingsong Wang, Fangyang Zheng}
\date{Source: arXiv:2608.00457v2 (CC BY 4.0)}
\begin{document}
\maketitle

\noindent\emph{Source and licence.} Balanced Bismut torsion-parallel fourfolds with constant holomorphic sectional curvature. Version arXiv:2608.00457v2 (CC BY 4.0), distributed by arXiv under the Creative Commons Attribution 4.0 International licence, http://creativecommons.org/licenses/by/4.0/, as recorded in the arXiv licence metadata for this exact version and shown on the abstract page https://arxiv.org/abs/2608.00457v2. Full licence notice: \texttt{licenses/arxiv-2608.00457-CC-BY-4.0.txt}.\par\smallskip

\noindent\textit{Editorial scope.} Counted unit: the article's proposition proprank2 (source0.tex lines 326--328). The article's complete original proof of it is delivered with it (source0.tex lines 330--387, one \texttt{proof} environment of 58 lines). No mathematics is written, repaired or re-derived: the statement and the proof are the article's own lines, in the article's own notation, and no retained line is substituted or re-flowed.  Retained uncounted as the source-local context the proof needs: lines 225--227; lines 247--253.  Nothing was omitted from any retained span, so the package carries no omission record.  No retained line contains a citation, so the wrapper carries no bibliography and no bibliography entry.  No retained line contains a figure environment, an image-inclusion command or drawing code, so no image is delivered and the package declares no asset dependency.  Editorial formatting only, in addition: an AMS \texttt{amsart} preamble that restates the article's own declarations and macros used by the retained lines, namely the environments \texttt{lemma}, \texttt{proposition} on separate counters, together with the standard packages those lines need.  The retained environments print in this document as Lemma 1, Proposition 1 in order of appearance, whereas the article prints the same items with the numbering of its own full document.  The complete native source of the article as distributed in the arXiv e-print for this exact version is bundled unchanged as \texttt{sources/206544/source0.tex}.
\begin{equation} \label{eq:Rb}
R^b_{i\bar{j}k\bar{\ell}} = \frac{c}{2} \big( \delta_{ij}\delta_{k\ell} + \delta_{i\ell}\delta_{kj} \big) - \frac{1}{2}\sum_{r=1}^n T^r_{ik} \overline{ T^r_{j\ell }} - \frac{3}{4}\sum_{r=1}^n \big( T^j_{ir} \overline{ T^k_{\ell r}} + T^{\ell}_{kr} \overline{ T^i_{j r}}  \big) + \frac{1}{4}\sum_{r=1}^n \big( T^{\ell}_{ir} \overline{ T^k_{j r}} + T^{j}_{kr} \overline{ T^i_{\ell r}}  \big),
\end{equation}

\begin{lemma} \label{lemmakernelB}
Let $(M^n,g)$ be a BTP manifold. Then, under any local unitary frame $e$, we have
\begin{equation} \label{eq:kernelB}
 R^b_{i\bar{j}k\bar{\ell}} = 0,  \ \ \ \ \ \ \ \mbox{if} \ e_k\in (\mbox{ker}(B))^{\perp} \  \mbox{and} \ \,e_{\ell} \in \mbox{ker}(B),
\end{equation}
for any $1\leq i,j\leq n$.
\end{lemma}

\begin{proposition} \label{proprank2}
Let $(M^4,g)$ be a balanced BTP manifold of dimension $4$ with constant Chern holomorphic sectional curvature. Then $r_B\neq 2$.
\end{proposition}

\begin{proof}
Assume, for a contradiction, that $r_B=2$. First let us choose a local unitary frame $e$ so that $\mbox{ker}(B)$ is spanned by $\{ e_3, e_4\}$. Then $T^3_{\ast \ast} = T^4_{\ast \ast}=0$. By the balanced assumption, $T^1_{12}=-T^3_{32}-T^4_{42}=0$. Similarly, $T^2_{12}=0$. Hence the only possibly non-zero torsion components under $e$ are the following:
\begin{equation} \label{eq:rank2}
E = (T^j_{i3})_{1\leq i,j\leq 2} = \left[ \begin{array}{cc} a & r\\p& -a \end{array} \right] , \ F = (T^j_{i4})_{1\leq i,j\leq 2} = \left[ \begin{array}{cc} b & s\\q& -b \end{array} \right],\ T^1_{34}=u, \ T^2_{34}=v.
\end{equation}
The traces of $E$ and $F$ are zero because the metric is balanced. By Lemma \ref{lemmakernelB}, we know that
\begin{equation} \label{eq:rank2a}
R^{b}_{i\bar{j}k\bar{\alpha}} =0,  \ \ \ \ \ \ \ \ \ \ \forall \ k\in \{ 1,2\}, \ \ \forall \ \alpha \in \{ 3,4\}, \ \  \forall \ 1\leq i,j \leq 4.
\end{equation}
On the other hand, since $g$ is BTP with constant Chern holomorphic sectional curvature, if we set \(\ell=\alpha\) in the equation (\ref{eq:Rb}), where \(k\in\{1,2\}\) and \(\alpha\in\{3,4\}\), then we get
\begin{equation} \label{eq:rank2b}
R^{b}_{i\bar{j}k\bar{\alpha}} = \frac{c}{2} \delta_{i\alpha} \delta_{kj} -\frac{1}{2} \sum_{r=1}^2 T^r_{ik}\overline{T^r_{j\alpha}}  + \frac{1}{4} \sum_{r=1}^4 \big( T^j_{kr}\overline{T^i_{\alpha r}} - 3 T^j_{ir}\overline{T^k_{\alpha r}} \big).
\end{equation}
When $j\in \{ 3,4\}$, the above equations (\ref{eq:rank2a}) and (\ref{eq:rank2b}) give us $\sum_r T^r_{ik} \overline{T^r_{34}} =0$. When $i\in \{ 3,4\}$, this leads to
\begin{equation}  \label{eq:rank2-1}
 E \left[ \begin{array}{c} \overline{u} \\ \overline{v} \end{array} \right] \ = \  F \left[ \begin{array}{c} \overline{u} \\ \overline{v} \end{array} \right]  = 0.
\end{equation}
Now let us assume that $j\in \{ 1,2\}$. When $i=k$, (\ref{eq:rank2a}) and (\ref{eq:rank2b}) give us
\begin{equation}  \label{eq:rank2-2}
 E_{1r} \overline{u} = F_{1r} \overline{u} = E_{2r} \overline{v} = F_{2r} \overline{v} = 0, \ \ \ \ \ \ \ \ \ \forall \ 1\leq r\leq 2.
\end{equation}
When $i,j \in \{ 1,2\}$ but $i\neq k$, we get from (\ref{eq:rank2a}) and (\ref{eq:rank2b}) that
$$  T^j_{k\beta} \overline{T^i_{34}} = 3  T^j_{i\beta} \overline{T^k_{34}}, \ \ \ \ \ \ \ \ \ \ \ j\in \{ 1,2\}, \ \beta \in \{ 3,4\}, \ \{i,k\} =\{ 1,2\}.$$
Taking $\beta=3$ and $(i,k)=(1,2)$ or $(2,1)$, we obtain $E_{2j}\overline{u}= 3E_{1j}\overline{v}$ and $E_{1j}\overline{v}= 3E_{2j}\overline{u}$, hence $E_{2j}\overline{u}= E_{1j}\overline{v}=0$. Similarly, by taking $\beta =4$ we get $F_{2j}\overline{u}= F_{1j}\overline{v}=0$. Thus we have
\begin{equation}  \label{eq:rank2-3}
 E_{2r} \overline{u} = E_{1r} \overline{v}=F_{2r} \overline{u} = F_{1r} \overline{v}=0, \ \ \ \ \ \ \ \ \ \forall \ 1\leq r\leq 2.
\end{equation}

We claim that $u=v=0$. To see this, assume on the contrary that $u\neq 0$ or $v\neq 0$. Then by (\ref{eq:rank2-2}) and (\ref{eq:rank2-3}) we have $E=F=0$. In this case the only possibly non-zero torsion components are $T^1_{34}=u$ and $T^2_{34}=v$. A non-trivial linear combination $w$ of $e_1$ and $e_2$ will then satisfy $T^w_{34}=0$ hence $T^w_{\ast \ast }=0$,  leading to $w\in \mbox{ker}(B)$, which is a contradiction to the assumption that $r_B=2$. So we know that we must have $u=v=0$.

Now let us come back to (\ref{eq:rank2a}) and (\ref{eq:rank2b}). Taking $j\in\{1,2\}$ and $i=\alpha$, we obtain
\begin{equation}  \label{eq:rank2-4}
 2cI = 3 E^{\ast}E - 2 E E^{\ast}, \ \ \ \ \ 2cI = 3 F^{\ast}F - 2 F F^{\ast} .
\end{equation}
Similarly, if we set $j\in \{ 1,2\}$ and $i\in \{ 3,4\}$ but $i\neq \alpha$, then we have
\begin{equation}  \label{eq:rank2-5}
  3 E^{\ast}F = 2 F E^{\ast}, \ \ \ \ \ 3 F^{\ast}E = 2 E F^{\ast} .
\end{equation}
Finally, by taking traces in (\ref{eq:rank2-4}), we get
\begin{equation}  \label{eq:rank2tr}
 4c \ = \, \ \parallel\! E\! \parallel^2 \ \, = \, \ \parallel\! F\! \parallel^2 .
\end{equation}

Let us proceed with the proof of Proposition \ref{proprank2}. We claim that $c$ must be $0$. Assume on the contrary that $c>0$. By writing out the elements of $E$ and $F$ in the equations (\ref{eq:rank2-4}) and (\ref{eq:rank2-5}), we obtain
\begin{eqnarray}
 && |p|^2=|r|^2 = 2c - |a|^2, \ \ \ \ \ \ |q|^2=|s|^2 = 2c-|b|^2, \label{eq:a}\\
 && a \bar{p} = \bar{a}r, \ \ \ \ \ \ b\bar{q} = \bar{b}s, \label{eq:b} \\
 && p\bar{q} = r\bar{s} = -a\bar{b}, \ \ \ \ \ \ a\bar{q}=\bar{b}r, \ \ \ \ \ \ a\bar{s}=\bar{b}p.   \label{eq:c}
\end{eqnarray}

Consider the row vectors $(a,p)$ and $(b,q)$ in ${\mathbb C}^2$. Equation (\ref{eq:a}) says that they have length $\sqrt{2c}>0$. These two vectors are perpendicular to each other because $a\bar{b}+p\bar{q}=0$ from the first equality of (\ref{eq:c}). Therefore these two vectors cannot be proportional to each other, namely, $aq-bp\neq 0$.

If $a=0$, then by the first equality of (\ref{eq:a}) we have $pr\neq 0$. From the first equality in (\ref{eq:c}), we then get $q=0$.  So by the second equality in (\ref{eq:a}) we know that $b\neq 0$. Now the  second equality of (\ref{eq:c}), $a\bar{q}=\bar{b}r$, would give us a contradiction.

If $a\neq 0$, then the first equality of (\ref{eq:b}) gives us $r=(a\bar{p})/\bar{a}$. Substituting this into the second equality of (\ref{eq:c}) gives $aq-bp=0$, a contradiction.

The above discussion indicates that the assumption $c>0$ always leads to a contradiction, so we must have $c=0$. By (\ref{eq:rank2tr}), we conclude that  $E=F=0$, hence $T=0$, violating the $r_B=2$ assumption. So for $n=4$ the $r_B=2$ case cannot occur. This completes the proof of the proposition.
\end{proof}

\end{document}
